% !TeX root = ../main.tex

\chapter{简介}

\section{一级节标题：介绍}

本模板\pkg{xynuthesis} 基于中国科学技术大学学位论文 LaTeX 模板（ustcthesis v4.0.0-beta.10）与清华大学学位论文 LaTeX 模板（ThuThesis）修改而成，
并依据《信阳师范学院研究生学位论文规范（2021年2月）》
2026年还在使用这一规范，以下简称《论文规范》）及论文示例（适用于专业硕士）的排版需求进行了专门适配。


由于官方提供的论文示例模板与《论文规范》描述并不一致，且《论文规范》中部分要求较为模糊，
因此在实现过程中参考了多个硕士毕业论文的格式并与《论文规范》和论文示例结合进行调整与取舍。


本模板适用于信阳师范大学专业硕士学位论文的撰写；学术硕士学位论文的撰写请参照《论文规范》及论文示例（适用于学术型硕士）。

本模板为非官方模板，并非学校或学院发布的正式 LaTeX 模板，无法保证与学校后续更新的排版规范完全同步。使用者应以《论文规范》及学院/研究生院的最新要求为准，并自行核对。

此外，本模板中的中英文封面与学位论文原创性声明页等内容，由于学校官方并没有提供明确的布局排版格式，所以本模板是参考学校提供的Word 示例，
在反复对照与逐步修改中实现版式逼近；因此其目标主要是视觉效果上的相似，在细节方面可能仍与 Word 版式存在差异。

\subsection{二级节标题：排版测试}
测试 XYNU Thesis LaTeX 模板的排版效果（比如检查字体、行距、页边距是否正常），
在 LaTeX 中不需要手动复制粘贴文字，
可以使用宏包自动生成。
测试英文排版（使用 lipsum 宏包）
测试中文排版（使用 zhlipsum 宏包）
\subsubsection{三级节标题：中文排版}

% 生成第 1 到第 3 段中文假文
\zhlipsum[1-2]

\subsubsection{三级节标题：英文排版}
% 生成第 1 到第 2 段英文假文
\lipsum[1-2]

\section{脚注}

Lorem ipsum dolor sit amet, consectetur adipiscing elit, sed do eiusmod tempor
incididunt ut labore et dolore magna aliqua.\footnote{Ut enim ad minim veniam, quis nostrud exercitation ullamco laboris
  nisi ut aliquip ex ea commodo consequat.
  Duis aute irure dolor in reprehenderit in voluptate velit esse cillum dolore
  eu fugiat nulla pariatur.}
